\documentclass[aps,prd,reprint,superscriptaddress,showpacs,preprintnumbers,nofootinbib,longbibliography,floatfix]{revtex4-2}
\usepackage[utf8]{inputenc}
\usepackage{amsmath,amssymb}
\usepackage{graphicx}
\PassOptionsToPackage{hyphens}{url}
\usepackage[colorlinks=true,allcolors=blue]{hyperref}
\usepackage{url}
\makeatletter
\g@addto@macro\UrlBreaks{\do\/\do\_\do\-\do\.}
\makeatother
\usepackage{booktabs}
\usepackage{ragged2e}
\newcommand{\fnl}{f_{\rm NL}}
\newcommand{\Pzeta}{\mathcal{P}_\zeta}
\newcommand{\eg}{e.g.}
\newcommand{\ie}{i.e.}
\newcommand{\etal}{\textit{et~al.}}

\newcommand{\paperVersion}{v2L.0.1}
\newcommand{\paperTimestamp}{September 2, 2026}

\begin{document}

\preprint{\paperVersion}

\title{The exact local non-Gaussianity of a matter-dominated contraction: \\
$f_{\rm NL}=-35/16$ and its orientation dependence}

\author{Houston Golden}
\email{houston@hubify.com}
\altaffiliation[ORCID: ]{\href{https://orcid.org/0009-0008-5616-5994}{0009-0008-5616-5994}}
\affiliation{Independent Researcher, Los Angeles, California, USA}

\date{\paperTimestamp}

\begin{abstract}
A matter-dominated cosmological contraction ($w=0$, $\epsilon=3/2$) generates local
primordial non-Gaussianity through the cubic self-interaction of the comoving-gauge
curvature perturbation $\zeta$. We report a from-scratch in-in computation of every
cubic vertex in Maldacena's action for this background, validated against the de~Sitter
and ultra-slow-roll limits before use. The isoceles squeezed-limit amplitude is
$f_{\rm NL}^{\rm local}=-35/16$, confirming the value obtained by Li, Quintin, Wang and
Cai (2016) at unit sound speed. We trace the long-standing discrepancy with $-35/8$,
originally reported by Cai, Xue, Brandenberger and Zhang (2009), to a uniform factor of
two in that work's amplitude-normalization step; the underlying shape function is
correct term by term. We report a new result: the squeezed limit is orientation
dependent, $f_{\rm NL}(\mu)=-35/16+(15/16)\mu^2$ with $\mu$ the cosine between the long
and short wavevectors (monopole $-15/8$, quadrupole $15/16$), sourced by the shear of
the non-attractor growing mode and analogous to the consistency-relation violation
identified by Namjoo, Firouzjahi and Sasaki in ultra-slow-roll inflation. A separate
$\delta N$ calculation on uniform-density slices returns $(5\epsilon-35)/8=-55/16$; this
is a different physical quantity, not a discrepancy, reconciled by $\zeta_\rho=2\zeta_c$
at linear order in a non-attractor phase. Survival of this amplitude through an explicit
nonsingular bounce is bounded, not resolved: linear transfer across three bounce
backgrounds obeys a universal suppression $0<T_{\fnl}\le1/2$, scheme-dependent at the
$1.6\times$ level, while the bounce's own cubic term -- found in the loop-quantum-cosmology
literature to enhance non-Gaussianity by orders of magnitude -- is not computed here. Only
a survey-scale bispectrum measurement can currently distinguish $-35/16$ from $-35/8$:
both fit DESI DR1 ($0.16\sigma$), while SPHEREx reaches $2.6$--$3.7\sigma$ between them.
\end{abstract}

\maketitle

\section{Setup}

We work in conformal time $\eta<0$ with $M_{\rm Pl}=1$, unit sound speed, and a matter
contraction $a(\eta)=c\,\eta^2$, for which the Hubble slow-roll parameter is fixed,
$\epsilon\equiv1-\mathcal H'/\mathcal H^2=3/2$, with $\mathcal H\equiv a'/a<0$. The
quadratic action for the comoving-gauge curvature perturbation $\zeta$
($h_{ij}=a^2e^{2\zeta}\delta_{ij}$, uniform-field gauge $\delta\phi=0$) is
$S_2=\int d\eta\,d^3x\,a^2\epsilon\,[\zeta'^2-(\partial\zeta)^2]$, with Bunch--Davies mode
functions fixed by the Wronskian normalization (not assumed). On super-Hubble scales the
physical mode is the growing, non-attractor solution $\zeta'=-3\zeta/\eta$
(equivalently $\dot\zeta=-\tfrac32H\zeta$), the analogue of the growing mode used by
Cai \etal~\cite{Cai:2009fn}.

We define the local bispectrum amplitude in the isoceles squeezed configuration
$k_1\ll k_2=k_3\equiv k$ exactly as in Refs.~\cite{Cai:2009fn,Li:2016}:
\begin{equation}
f_{\rm NL}^{\rm local}\equiv\lim_{k_1\to0}\frac{10}{3}\,\frac{\mathcal A(k_1,k,k)}{2k^3},
\end{equation}
where $\langle\zeta_{\mathbf k_1}\zeta_{\mathbf k_2}\zeta_{\mathbf k_3}\rangle
=(2\pi)^7\delta^3(\textstyle\sum\mathbf k_i)\,\Pzeta^2\,\mathcal A/\prod k_i^3$, equivalent
to the local template $\zeta=\zeta_g+\tfrac35\fnl\zeta_g^2$. This limit is independent of
the end-of-contraction time $\eta_B$ at leading order ($\mathcal O(k^2\eta_B^2)$
corrections) but, as we show below, is not independent of the orientation of the short
pair relative to the long mode.

We use the cubic action derived by Maldacena~\cite{Maldacena:2002vr} in comoving gauge
for a canonical field, retaining every term at $\eta_{\rm sr}=0$: two bulk pieces
proportional to $\epsilon^2$, one boundary/field-redefinition term, and the two
$\epsilon^3$ pieces that involve $\tilde\chi\equiv\partial^{-2}\zeta'$. All four bulk
vertices and the field redefinition are evaluated by first-order in-in perturbation
theory, $\langle\zeta^3(\eta_*)\rangle=-2\,{\rm Im}\int^{\eta_*}\!d\eta\,\langle
\zeta^3(\eta_*)L_{\rm int}(\eta)\rangle$, with every Wick contraction summed explicitly
(no symmetry factor inserted by hand) and every time integral performed exactly before
taking $|k\eta_B|\to0$.

\section{Validation}

Before applying this pipeline to the matter contraction we reproduce two independent
literature results with the same code: (i) the de~Sitter, constant-$\epsilon$ bispectrum
of Ref.~\cite{Maldacena:2002vr}, matched term by term; and (ii) the ultra-slow-roll
result of Namjoo, Firouzjahi and Sasaki~\cite{Namjoo:2012}, $f_{\rm NL}=5/2$, which
depends on the same field-redefinition structure that we use below. Both match exactly.

\section{Result}

Table~\ref{tab:vertices} gives the squeezed-limit contribution of each vertex, from
scratch, at $\epsilon=3/2$. The total reproduces the isoceles squeezed value
\begin{equation}
\boxed{\,f_{\rm NL}^{\rm local}\big|_{k_1\ll k_2=k_3}=-\frac{35}{16}\,}
\end{equation}
and, at general orientation $\mu\equiv\hat{\mathbf k}_1\!\cdot\!\hat{\mathbf k}$,
\begin{equation}
f_{\rm NL}^{\rm sq}(\mu)=-\frac{35}{16}+\frac{15}{16}\mu^2,
\label{eq:orientation}
\end{equation}
i.e.\ a monopole $-15/8$ and a quadrupole $15/16$ under angular averaging. The
equilateral and folded values are $-255/128$ and $-9/8$, matching
Ref.~\cite{Li:2016}'s $c_s=1$ limit exactly.

\begin{table}[t]
\caption{Squeezed-limit contribution of each cubic vertex at $\epsilon=3/2$, from a
from-scratch in-in computation. $\mu$-dependence is shown where nonzero.}
\label{tab:vertices}
\begin{ruledtabular}
\begin{tabular}{lc}
vertex & $f^{\rm sq}(\mu)$ \\
\hline
field redefinition (local + non-local part) & $-\dfrac{25}{16}-\dfrac{15}{16}\mu^2$ \\
$\zeta\zeta'^2\ (\epsilon^2-\epsilon^3/2)$ & $-\dfrac{5}{32}$ \\
$\zeta(\partial\zeta)^2$ & $0$ \\
$\zeta'\partial\zeta\!\cdot\!\partial\tilde\chi$ & $\dfrac{15}{8}\mu^2$ \\
$\zeta(\partial_i\partial_j\tilde\chi)^2$ & $-\dfrac{15}{32}$ \\
\hline
\textbf{total} & $-\dfrac{35}{16}+\dfrac{15}{16}\mu^2$ \\
\end{tabular}
\end{ruledtabular}
\end{table}

\textit{Origin of the discrepancy with Cai \etal~\cite{Cai:2009fn}.} Reading their
printed shape function (their Eq.~37) as the sum of three distinct $(5,2,2)$-type
monomials, it is identical to our vertex sum, and to Li \etal's Eq.~4.19 at $c_s=1$
(difference zero, checked symbolically). Cai \etal's three quoted amplitudes -- local
$-35/8$, equilateral $-255/64$, folded $-9/4$ -- are each exactly twice the values in
this Letter, in all three configurations. The origin is a uniform factor of two carried
into their amplitude-normalization step, not an error in any vertex, commutator, or Wick
contraction. Li \etal~\cite{Li:2016} inherit Cai \etal's correct polynomial and are
therefore not an independent check; Quintin \etal~\cite{Quintin:2015} cite $-35/16$
without an independent computation. The present result is, to our knowledge, the first
computation of this amplitude that does not transcribe a per-vertex expression from
either.

\textit{Orientation dependence and its origin.} Equation~(\ref{eq:orientation}) arises
because the growing, non-attractor mode is not locally an FRW patch as $k_1\to0$: its
extrinsic-curvature perturbation carries both a trace, $\delta K\to\epsilon\dot\zeta$,
and a shear of the same order, $\sigma^i{}_j=(\hat k_i\hat k_j-\tfrac13\delta_{ij})\,
\epsilon\dot\zeta$, sourced by the comoving-gauge shift $\partial\tilde\chi_L\propto
\hat{\mathbf k}_1\zeta_L'/k_1$. The $\mu^2$ term comes entirely from the two vertices in
which the long mode enters through $\partial\tilde\chi$. This is the matter-contraction
analogue of the consistency-relation violation Namjoo, Firouzjahi and
Sasaki~\cite{Namjoo:2012} found in ultra-slow-roll inflation, and is consistent with the
general non-attractor clock argument of Chen, Namjoo and Wang~\cite{Chen:2013aj}: a
squeezed-limit amplitude that survives the $k_1\to0$ limit is the generic signature of a
non-attractor background, not an artifact of this one.

\textit{Cross-check by a different route.} A separate-universe ($\delta N$) calculation
on the same background but evaluated on \emph{uniform-density} slices gives, for general
$\epsilon$, $f_{\rm NL}^\rho=(5\epsilon-7)\cdot 5/8$ reduced at $\epsilon=3/2$ to
$(5\epsilon-35)/8=-55/16$, while the comoving-slice $\delta N$ (an isotropic, shear-free
approximation) gives $f_{\rm NL}^c=-5$ for all $\epsilon$. Neither equals the in-in
result above, and this is expected rather than alarming: at linear order
$\zeta_\rho=2\zeta_c$ in a non-attractor phase, so uniform-density and comoving slices
differ already at $\mathcal O(1)$; and the isotropic separate-universe approximation
discards the shear that produces Eq.~(\ref{eq:orientation}) by construction, keeping
only the trace. The $\delta N$ values are correct for their own variables and are not
forced to agree with the in-in result; they instead confirm that the gap between $-35/16$
and Cai \etal's $-35/8$ is a clean factor of two rather than a definitional mismatch, and
they leave a single well-posed open calculation -- an anisotropic (Bianchi-I)
separate-universe treatment of the long mode's shear response -- as the remaining
independent check of the monopole and quadrupole in Eq.~(\ref{eq:orientation}).

\section{Transmission through the bounce: stated honestly}

Equation~(\ref{eq:orientation}) is a statement about the contraction phase alone. Whether
it survives into an observable post-bounce spectrum depends on physics this Letter does
not compute in full, and we state the current evidence precisely. Linear mode-mixing
across three explicit nonsingular bounce backgrounds -- an LQC effective quasi-dust
bounce, an analytic non-LQC bounce, and the Quintin \etal~\cite{Quintin:2015}-type
parametrization -- gives a transmission coefficient
$T_{\fnl}=[1-\rho(\eta_h)]/2$ for the growing-mode-dominated adiabatic vacuum, obeying a
universal bound $0<T_{\fnl}\le1/2$: linear transfer can only suppress the amplitude, by
scheme- and background-dependent factors in the range $0.165$--$0.409$ (a $1.64\times$
spread between two mode-function conventions on the same background). This supersedes an
earlier, narrower statement about the transfer of the sub-dominant constant branch alone.
What this Letter does not compute is the bounce's own cubic self-interaction term, which
the loop-quantum-cosmology literature finds can \emph{enhance} non-Gaussianity by orders
of magnitude through gravitational self-interactions in the third-order Hamiltonian, even
in the absence of a potential~\cite{AgulloBollietSreenath:2017}. Any assumption that the
bounce transmits $\fnl$ with a coefficient near unity is therefore not supported by
either calculation reported here.

\section{Observational reach}

Table~\ref{tab:reach} compares the two candidate amplitudes against current and
near-future large-scale-structure constraints, using the scale-dependent-bias signature
of local non-Gaussianity~\cite{Dalal:2007cu}. DESI DR1 cannot currently distinguish
$-35/16$ from $-35/8$; SPHEREx is the first instrument projected to.

\begin{table}[t]
\caption{Discriminating power of current and near-future bispectrum/scale-dependent-bias
constraints on $f_{\rm NL}^{\rm local}=-35/16$ vs.\ $-35/8$.}
\label{tab:reach}
\begin{ruledtabular}
\small
\begin{tabular}{lccc}
survey & $\sigma_{f_{\rm NL}}$ & tension, $-35/16$ & separation \\
\hline
DESI DR1 & -- & $0.16\sigma$ & $\lesssim0.2\sigma$ \\
SPHEREx, bispectrum & $0.7$ & -- & $2.6$--$3.1\sigma$ \\
SPHEREx, target ($P{+}B$) & $0.5$ & -- & $3.7$--$4.4\sigma$ \\
\end{tabular}
\end{ruledtabular}
\end{table}

SPHEREx sensitivities are quoted from Ref.~\cite{Heinrich:2023}
(\S~2311.13082, abstract) and Ref.~\cite{Dore:2014}; DESI DR1 comparison uses the
public headline scale-dependent-bias constraint. Full construction of both channels,
plus a Press--Schechter-level primordial-black-hole comparison against
Ref.~\cite{Choudhury:2025} and a PTA free-spectrum consistency check against the
induced-gravitational-wave slope of Ref.~\cite{Papanikolaou:2025}, are given in
Ref.~\cite{Golden2026Multichannel}.

\section{Summary}

A from-scratch in-in computation, validated against two independent literature limits,
gives $f_{\rm NL}^{\rm local}=-35/16$ for a canonical matter-dominated contraction, with
a new orientation-dependent squeezed limit $-35/16+(15/16)\mu^2$ sourced by the
non-attractor shear. The long-discussed factor-of-two discrepancy with
Ref.~\cite{Cai:2009fn} is traced to a single normalization step, not to the underlying
shape function, which is correct. What remains open is not the amplitude at the end of
the contraction but its fate through the bounce: the linear transfer is bounded,
$T_{\fnl}\le1/2$, and the bounce's intrinsic nonlinear contribution is not yet computed.
Only a survey-scale bispectrum measurement, at the sensitivity SPHEREx targets, can
currently discriminate the corrected amplitude from the value it replaces.

\section*{Data and Code Availability}
The from-scratch symbolic derivation, its two literature validation checks, the
separate-universe cross-check, the bounce-transmission computation, and the
multi-channel reach calculation are committed at
\url{https://github.com/Hubify-Projects/bigbounce} under
\texttt{research/theory\_audit/}, \texttt{research/cubic\_bounce\_transmission/}, and
\texttt{research/track\_a3\_multichannel/} (commits referenced in
Refs.~\cite{Golden2026TheoryAudit,Golden2026Transmission,Golden2026Multichannel}). The
parent forecast paper and its own code/data are permanently archived at Zenodo,
DOI~10.5281/zenodo.21461881~\cite{Golden2026P2Zenodo}.

\section*{AI Usage Disclosure}
Symbolic derivations reported here were performed by deterministic sympy scripts written
and executed under the author's direction; every intermediate result is printed by the
committed scripts and independently checked against the cited literature limits before
use. Claude (Anthropic, Opus-tier) assisted with literature cross-checks, prose drafting,
and LaTeX formatting under the author's direction; all physics content, equations, and
numerical results were verified by the author against the committed script output.

\bibliography{references}

\end{document}
